% notes-l-structures.tex
% Topic: Languages and Structures

\begin{remark*}[Exposition]
Following Kirby's model-theoretic convention, a language is the formal
vocabulary in which structures are described. A structure then interprets that
vocabulary on a concrete universe.
\end{remark*}

\begin{definitionbox}{Definition (First-Order Language)}
\begin{definition}[First-Order Language]\label{def:first-order-language}
A \textbf{first-order language} $\mathcal{L}$ consists of a set
$\mathcal{R}_{\mathcal{L}}$ of relation symbols, a set
$\mathcal{F}_{\mathcal{L}}$ of function symbols, a set
$\mathcal{C}_{\mathcal{L}}$ of constant symbols, and an arity assignment for
each relation and function symbol.
\end{definition}
\end{definitionbox}

\begin{remark*}[Standard quantified statement]
$\mathcal{L}=(\mathcal{R}_{\mathcal{L}},\mathcal{F}_{\mathcal{L}},
\mathcal{C}_{\mathcal{L}},\operatorname{ar})$, where
$\operatorname{ar}(R)\in\mathbb{N}_{\geq 1}$ for each
$R\in\mathcal{R}_{\mathcal{L}}$ and
$\operatorname{ar}(f)\in\mathbb{N}_{\geq 1}$ for each
$f\in\mathcal{F}_{\mathcal{L}}$.
\end{remark*}

\begin{remark*}[Interpretation]
The language is also called a vocabulary or a signature. It specifies which
names may occur in formulas and what arity each non-logical symbol has, but it
does not yet say what those symbols denote.
\end{remark*}

\DefinitionalRoot

\begin{definitionbox}{Definition ($\mathcal{L}$-Structure)}
\begin{definition}[$\mathcal{L}$-Structure]\label{def:l-structure}
Let $\mathcal{L}$ be a first-order language. An
\textbf{$\mathcal{L}$-structure} $\mathcal{M}$ consists of a nonempty set
$M$, called the universe of $\mathcal{M}$, together with interpretations of
all symbols of $\mathcal{L}$ on $M$: each relation symbol $R$ of arity $n$ is
assigned a relation $R^{\mathcal{M}}\subseteq M^n$, each function symbol $f$
of arity $n$ is assigned a function $f^{\mathcal{M}}\colon M^n\to M$, and
each constant symbol $c$ is assigned an element $c^{\mathcal{M}}\in M$.
\end{definition}
\end{definitionbox}

\begin{remark*}[Standard quantified statement]
$\mathcal{M}$ is an $\mathcal{L}$-structure iff
$M\neq\emptyset$, $R^{\mathcal{M}}\subseteq M^{\operatorname{ar}(R)}$ for
each relation symbol $R$, $f^{\mathcal{M}}\colon
M^{\operatorname{ar}(f)}\to M$ for each function symbol $f$, and
$c^{\mathcal{M}}\in M$ for each constant symbol $c$.
\end{remark*}

\begin{remark*}[Interpretation]
A structure is a language made concrete. The language gives the names and
arities; the structure supplies the underlying set and the actual relations,
operations, and distinguished elements named by the language.
\end{remark*}

\begin{dependencies}
\begin{itemize}
  \item \hyperref[def:first-order-language]{First-Order Language}
\end{itemize}
\end{dependencies}

\begin{definitionbox}{Definition (First-Order Sublanguage)}
\begin{definition}[First-Order Sublanguage]\label{def:first-order-sublanguage}
Let $\mathcal{L}_0$ and $\mathcal{L}$ be first-order languages.
The language $\mathcal{L}_0$ is a \textbf{sublanguage} of $\mathcal{L}$ if
each relation, function, and constant symbol of $\mathcal{L}_0$ is a symbol of
$\mathcal{L}$ with the same arity.
\end{definition}
\end{definitionbox}

\begin{remark*}[Standard quantified statement]
$\mathcal{L}_0\subseteq\mathcal{L}$ iff
$\mathcal{R}_{\mathcal{L}_0}\subseteq\mathcal{R}_{\mathcal{L}}$,
$\mathcal{F}_{\mathcal{L}_0}\subseteq\mathcal{F}_{\mathcal{L}}$,
$\mathcal{C}_{\mathcal{L}_0}\subseteq\mathcal{C}_{\mathcal{L}}$, and all
arities agree.
\end{remark*}

\begin{remark*}[Interpretation]
A sublanguage forgets some vocabulary but does not change the meaning or arity
of the symbols it keeps.
\end{remark*}

\begin{dependencies}
\begin{itemize}
  \item \hyperref[def:first-order-language]{First-Order Language}
\end{itemize}
\end{dependencies}

\begin{definitionbox}{Definition ($\mathcal{L}$-Structure Reduct)}
\begin{definition}[$\mathcal{L}$-Structure Reduct]\label{def:l-structure-reduct}
Let $\mathcal{L}_0$ be a sublanguage of $\mathcal{L}$, and let
$\mathcal{M}$ be an $\mathcal{L}$-structure. The
\textbf{$\mathcal{L}_0$-reduct} of $\mathcal{M}$ is the
$\mathcal{L}_0$-structure $\mathcal{M}\!\upharpoonright\!\mathcal{L}_0$
with the same universe as $\mathcal{M}$ and with the interpretations of the
symbols of $\mathcal{L}_0$ inherited from $\mathcal{M}$.
\end{definition}
\end{definitionbox}

\begin{remark*}[Standard quantified statement]
If $\mathcal{L}_0\subseteq\mathcal{L}$ and $\mathcal{M}$ is an
$\mathcal{L}$-structure, then
$|\mathcal{M}\!\upharpoonright\!\mathcal{L}_0|=|\mathcal{M}|$, and every
$\mathcal{L}_0$-symbol has the same interpretation in
$\mathcal{M}\!\upharpoonright\!\mathcal{L}_0$ as in $\mathcal{M}$.
\end{remark*}

\begin{remark*}[Interpretation]
A reduct is obtained by forgetting structure. It keeps the same underlying
set, but only remembers the symbols belonging to the smaller language.
\end{remark*}

\begin{dependencies}
\begin{itemize}
  \item \hyperref[def:first-order-sublanguage]{First-Order Sublanguage}
  \item \hyperref[def:l-structure]{$\mathcal{L}$-Structure}
\end{itemize}
\end{dependencies}

\begin{definitionbox}{Definition ($\mathcal{L}$-Structure Expansion)}
\begin{definition}[$\mathcal{L}$-Structure Expansion]\label{def:l-structure-expansion}
Let $\mathcal{L}_0$ be a sublanguage of $\mathcal{L}$. If
$\mathcal{A}$ is an $\mathcal{L}_0$-structure and $\mathcal{M}$ is an
$\mathcal{L}$-structure whose $\mathcal{L}_0$-reduct is $\mathcal{A}$, then
$\mathcal{M}$ is an \textbf{expansion} of $\mathcal{A}$ to $\mathcal{L}$.
\end{definition}
\end{definitionbox}

\begin{remark*}[Standard quantified statement]
$\mathcal{M}$ expands $\mathcal{A}$ from $\mathcal{L}_0$ to $\mathcal{L}$ iff
$\mathcal{M}\!\upharpoonright\!\mathcal{L}_0=\mathcal{A}$.
\end{remark*}

\begin{remark*}[Interpretation]
An expansion adds interpretations for new symbols while preserving all
previously interpreted vocabulary.
\end{remark*}

\begin{dependencies}
\begin{itemize}
  \item \hyperref[def:first-order-sublanguage]{First-Order Sublanguage}
  \item \hyperref[def:l-structure-reduct]{$\mathcal{L}$-Structure Reduct}
\end{itemize}
\end{dependencies}

\begin{definitionbox}{Definition (Homomorphism of $\mathcal{L}$-Structures)}
\begin{definition}[Homomorphism of $\mathcal{L}$-Structures]\label{def:l-structure-homomorphism}
Let $\mathcal{M}$ and $\mathcal{N}$ be $\mathcal{L}$-structures. A
\textbf{homomorphism} $h\colon\mathcal{M}\to\mathcal{N}$ is a function
$h\colon M\to N$ that preserves constants, functions, and relations of
$\mathcal{L}$.
\end{definition}
\end{definitionbox}

\begin{remark*}[Standard quantified statement]
$h\colon\mathcal{M}\to\mathcal{N}$ is a homomorphism iff
$h(c^{\mathcal{M}})=c^{\mathcal{N}}$ for every constant symbol $c$,
$h(f^{\mathcal{M}}(\bar a))=f^{\mathcal{N}}(h(\bar a))$ for every function
symbol $f$ and tuple $\bar a$, and whenever
$R^{\mathcal{M}}(\bar a)$ holds, $R^{\mathcal{N}}(h(\bar a))$ holds.
\end{remark*}

\begin{remark*}[Interpretation]
Homomorphisms preserve the positive algebraic and relational structure named
by the language.
\end{remark*}

\begin{dependencies}
\begin{itemize}
  \item \hyperref[def:l-structure]{$\mathcal{L}$-Structure}
\end{itemize}
\end{dependencies}

\begin{definitionbox}{Definition (Embedding of $\mathcal{L}$-Structures)}
\begin{definition}[Embedding of $\mathcal{L}$-Structures]\label{def:l-structure-embedding}
Let $\mathcal{M}$ and $\mathcal{N}$ be $\mathcal{L}$-structures. An
\textbf{embedding} $e\colon\mathcal{M}\to\mathcal{N}$ is an injective function
$e\colon M\to N$ that preserves and reflects every relation symbol, preserves
every function symbol, and preserves every constant symbol.
\end{definition}
\end{definitionbox}

\begin{remark*}[Standard quantified statement]
$e$ is an embedding iff $e$ is injective,
$e(c^{\mathcal{M}})=c^{\mathcal{N}}$,
$e(f^{\mathcal{M}}(\bar a))=f^{\mathcal{N}}(e(\bar a))$, and
$R^{\mathcal{M}}(\bar a)$ iff $R^{\mathcal{N}}(e(\bar a))$.
\end{remark*}

\begin{remark*}[Interpretation]
An embedding identifies $\mathcal{M}$ with a copy of itself inside
$\mathcal{N}$, preserving the truth of atomic relational information both
forwards and backwards.
\end{remark*}

\begin{dependencies}
\begin{itemize}
  \item \hyperref[def:l-structure]{L-Structure}
  \item \hyperref[def:l-structure-homomorphism]{Homomorphism of $\mathcal{L}$-Structures}
\end{itemize}
\end{dependencies}

\begin{definitionbox}{Definition (Isomorphism of $\mathcal{L}$-Structures)}
\begin{definition}[Isomorphism of $\mathcal{L}$-Structures]\label{def:l-structure-isomorphism}
An \textbf{isomorphism} of $\mathcal{L}$-structures is a bijective embedding.
\end{definition}
\end{definitionbox}

\begin{remark*}[Standard quantified statement]
$i\colon\mathcal{M}\to\mathcal{N}$ is an isomorphism iff $i$ is an embedding
and $i$ is bijective.
\end{remark*}

\begin{remark*}[Interpretation]
Isomorphic structures have the same $\mathcal{L}$-structure up to a relabeling
of their underlying elements.
\end{remark*}

\begin{dependencies}
\begin{itemize}
  \item \hyperref[def:l-structure-embedding]{Embedding of $\mathcal{L}$-Structures}
\end{itemize}
\end{dependencies}

\begin{definitionbox}{Definition ($\mathcal{L}$-Structure Automorphism)}
\begin{definition}[$\mathcal{L}$-Structure Automorphism]\label{def:l-structure-automorphism}
An \textbf{automorphism} of an $\mathcal{L}$-structure $\mathcal{M}$ is an
isomorphism from $\mathcal{M}$ to itself.
\end{definition}
\end{definitionbox}

\begin{remark*}[Standard quantified statement]
$\sigma$ is an automorphism of $\mathcal{M}$ iff
$\sigma\colon\mathcal{M}\to\mathcal{M}$ is an isomorphism.
\end{remark*}

\begin{remark*}[Interpretation]
Automorphisms are the symmetries of a structure. They permute the universe
while preserving every relation, function, and constant named by the language.
\end{remark*}

\begin{dependencies}
\begin{itemize}
  \item \hyperref[def:l-structure-isomorphism]{Isomorphism of $\mathcal{L}$-Structures}
\end{itemize}
\end{dependencies}
